\documentclass[11pt]{article}
\usepackage[top=0.65in, bottom=0.65in, left=0.75in, right=0.75in]{geometry}
\usepackage{enumitem}
\usepackage{hyperref}
\usepackage{xcolor}
\usepackage{titlesec}
\usepackage{parskip}
\usepackage{booktabs}
\newcommand{\uline}[1]{\underline{#1}}

\pagestyle{empty}

\definecolor{sectionblue}{RGB}{0, 51, 102}
\titleformat{\section}{\color{sectionblue}\large\bfseries}{}{0em}{}[\titlerule]
\titlespacing{\section}{0pt}{6pt}{4pt}
\titleformat{\subsection}[runin]{\bfseries}{}{0em}{}[.]

\hypersetup{
  colorlinks=true,
  linkcolor=sectionblue,
  urlcolor=sectionblue,
  pdftitle={Ishraq Tashdid --- Curriculum Vitae},
  pdfauthor={Ishraq Tashdid}
}

\setlength{\emergencystretch}{3em}
\setlength{\parskip}{2pt}
\setlength{\parsep}{0pt}
\renewcommand{\baselinestretch}{0.97}
\setlist[itemize]{noitemsep, topsep=1pt, left=0pt, parsep=0pt}

\begin{document}

% ─── HEADER ──────────────────────────────────────────────────────────────────
\begin{center}
    {\LARGE \textbf{ISHRAQ TASHDID}}\\
    \vspace{5pt}
    \href{mailto:tashdid.ishraq@gmail.com}{tashdid.ishraq@gmail.com} \;|\;
    \href{mailto:ishraq.tashdid@ucf.edu}{ishraq.tashdid@ucf.edu} \;|\; +1 (689) 286-7719\\
    Orlando, FL 32817 \;|\;
    \href{https://www.linkedin.com/in/ishraqtashdid}{LinkedIn} \;|\;
    \href{https://github.com/IshraqAtUCF}{GitHub} \;|\;
    \href{https://scholar.google.com/citations?user=gYcWmKYAAAAJ}{Google Scholar} \;|\;
    \href{https://ishraqtashdid.com}{ishraqtashdid.com}
\end{center}

\vspace{4pt}

% ─── RESEARCH INTERESTS ──────────────────────────────────────────────────────
\section*{Research Interests}
Automated vulnerability reasoning and threat modeling (CWE/CVE/CAPEC); LLM-assisted security
analysis across hardware--firmware--software boundaries; formal and property-based security
verification; chiplet and SiP authentication; secure SoC design and supply chain integrity.


% ─── EDUCATION ───────────────────────────────────────────────────────────────
\section*{Education}
\textbf{University of Central Florida (UCF)} \hfill \textit{Jan.\ 2024--Present}\\
Ph.D. in Computer Engineering (Joint PhD/MS) \hfill GPA: \textbf{3.96}/4.00\\
Advisor: Dr. Sazadur Rahman, Department of Electrical \& Computer Engineering\\
\href{https://www.ece.ucf.edu/doctoral-student-earns-ucf-trustee-fellowship/}{\uline{\textbf{Trustee's Doctoral Fellow}}} (4-year, \$25{,}000/year) \quad
\href{https://www.ece.ucf.edu/future-faculty-laureates-ffl-program}{\uline{\textbf{Future Faculty Laureate}}}, UCF ECE

\vspace{4pt}
\textbf{Bangladesh University of Engineering and Technology (BUET)} \hfill \textit{Feb.\ 2017--Jun.\ 2022}\\
B.Sc.\ in Electrical and Electronic Engineering \hfill CGPA: 3.48/4.00


% ─── RESEARCH EXPERIENCE ─────────────────────────────────────────────────────
\section*{Research Experience}

\textbf{Graduate Research Assistant} --- University of Central Florida \hfill \textit{Jan.\ 2024--Present}\\
\textit{Hardware Security \& Trustworthy Computing Lab, advised by Dr.\ Sazadur Rahman}

\begin{itemize}

  \item \textbf{LLM-driven threat modeling \& automated vulnerability reasoning (DAC~2026):}
  Built \emph{ATLAS}, a framework that ingests CWE/CVE/CAPEC vulnerability databases,
  auto-generates structured threat models for any SoC asset, and translates them into formal
  security properties verified by JasperGold. Detected $39/48$ known CWEs across three
  HACK@DAC benchmarks ($>$82\% correct property generation). Conducted in collaboration with
  \textbf{Jonathan Valamehr (Intel Research)}, whose co-authorship reflects Intel's institutional
  validation of the framework. Subsequently invited to present ATLAS to Intel's AI research
  community at the \textbf{Intel AI Forum (Feb.\ 2026)}. Framework is directly applicable to
  firmware and software vulnerability workflows that reason over the same threat taxonomies.

  \item \textbf{Automated security property generation \& cross-abstraction analysis:}
  Developed pipelines combining static analysis, AST-level structural reasoning, execution
  traces, and formal verification feedback to generate precise security assertions with
  minimal human effort, reducing reliance on manually crafted security rules at every
  abstraction layer from RTL through system firmware. This work forms the analytical
  backbone of both ATLAS and subsequent cross-abstraction security projects.

  \item \textbf{Chiplet authentication via PUF and MPC (HOST~2026):}
  Designed \emph{InterPUF}, a distributed authentication framework for reconfigurable
  System-in-Package (SiP) interposers. Embedded a differential delay PUF in the
  interconnect fabric and secured evaluation using multi-party computation (MPC), ensuring
  raw PUF signatures never leave the device. Achieved $<$0.23\% area and $<$0.072\% power
  overhead; ML modeling attacks converged at $\sim$47\% accuracy (random-guess level),
  demonstrating resistance to state-of-the-art ML-based reverse engineering. Addresses
  supply chain integrity concerns in multi-vendor chiplet ecosystems flagged by the
  CHIPS and Science Act (2022) and DoD supply chain security reports.

  \item \textbf{Routing-based PUF authentication (GOMACTech~2026):}
  Designed \emph{RoutePUF}, a routing-based physically unclonable function for secure
  interposer authentication in chiplet systems, extending the PUF-based authentication line
  of work to a government/defense-microelectronics audience via the Government Microcircuit
  Applications and Critical Technology Conference, a forum jointly sponsored by the U.S.
  Department of Defense.

  \item \textbf{Adaptive logic obfuscation for SoC IP protection (ICCD~2025):}
  Developed \emph{ECOLogic}, a framework enabling circular, obfuscated, and adaptive logic
  via eFPGA-augmented SoCs. The approach allows post-fabrication logic reconfiguration to
  resist reverse engineering and hardware Trojan insertion while maintaining functional
  correctness and low overhead.

  \item \textbf{Macro placement automation (MLCAD~2025 --- Best Paper Nominee):}
  Designed human-inspired reinforcement learning frameworks for chip floorplanning to
  improve post-route routability and reliability in heterogeneous SoCs. Demonstrated that
  incorporating human design intuition into RL reward modeling yields measurable gains over
  PPA-only optimization baselines.

  \item \textbf{Scalable chiplet trust architectures:}
  Developed \emph{SAFE-SiP} (GLSVLSI~2025) and \emph{AuthenTree} (PAINE 2025), two
  complementary MPC-based distributed trust frameworks addressing authentication scalability
  in chiplet-based heterogeneous integration. AuthenTree achieves logarithmic communication
  complexity in the number of chiplets, enabling practical deployment in large-scale SiP designs.

\end{itemize}


% ─── INDUSTRY EXPERIENCE ─────────────────────────────────────────────────────
\section*{Industry Experience}

\textbf{XPU Security Research Intern} --- \textbf{Intel Corporation}, Folsom, CA
\hfill \textit{May 2026--Present}
\begin{itemize}
  \item Competitive research internship in Intel's XPU (accelerator) security division.
  Selected following an invited talk at the Intel AI Forum on LLM-driven threat modeling
  for SoC security (Feb.\ 2026). Focused on security research for next-generation
  accelerator architectures, working with Intel's AI and hardware security research community.
  Supervisor: James Cavanaugh, Intel.
\end{itemize}

\vspace{4pt}
\textbf{Hardware Verification Engineer} --- \href{https://xcelerium.com/}{XCelerium}, Irvine, CA
{\small (via \href{https://www.dsinnovators.com/}{DSi}, Bangladesh)} \hfill \textit{Jan.\ 2022--Dec.\ 2023}

\begin{itemize}
  \item Verified custom \textbf{RISC-V} processor implementations at the ISA and
  \textbf{system level} using SystemVerilog, UVM, and cocotb, developing deep familiarity
  with the hardware--software boundary where firmware executes and privilege/mode
  transitions occur.

  \item Verified \textbf{AXI memory controllers} and \textbf{APB-AXI bridges},
  including CDC edge-cases and assertion-based checks on bus-level access control ---
  the same protocol surfaces targeted by firmware and software security assessments.

  \item Built \textbf{Python-driven automation frameworks} for coverage convergence,
  testbench orchestration, and results analysis, accelerating sign-off across multiple
  tapeout projects.
\end{itemize}


% ─── PUBLICATIONS ────────────────────────────────────────────────────────────
\section*{Publications}

\begin{itemize}[leftmargin=*]

  \item \textbf{Ishraq Tashdid}, Kimia Tasnia, Alexander Garcia, Jonathan Valamehr, Sazadur Rahman,
    \textit{``ATLAS: AI-Assisted Threat-to-Assertion Learning for System-on-Chip Security
    Verification,''} \textbf{DAC 2026} (63rd Design Automation Conference), Long Beach, CA,
    USA. \emph{Accepted.} \href{https://arxiv.org/abs/2603.01170}{arXiv:2603.01170}

  \item \textbf{Ishraq Tashdid},
    \textit{``RoutePUF: A Routing-Based Physically Unclonable Function for Secure Interposer
    Authentication in Chiplet Systems,''} \textbf{GOMACTech 2026} (Government Microcircuit
    Applications and Critical Technology Conference). \emph{Accepted.}

  \item \textbf{Ishraq Tashdid}, Tasnuva Farheen, Sazadur Rahman,
    \textit{``InterPUF: Distributed Authentication via Physically Unclonable Functions and
    Multi-party Computation for Reconfigurable Interposers,''} \textbf{HOST 2026}
    (IEEE Int'l Symposium on Hardware Oriented Security and Trust), Washington, D.C., USA.
    \emph{Accepted.} \href{https://arxiv.org/abs/2601.11368}{arXiv:2601.11368}

  \item \textbf{Ishraq Tashdid}, Dewan Saiham, Nafisa Anjum, Tasnuva Farheen, Sazadur Rahman,
    \textit{``ECOLogic: Enabling Circular, Obfuscated, and Adaptive Logic via
    eFPGA-Augmented SoCs,''} \textbf{ICCD 2025} (IEEE Int'l Conference on Computer Design),
    Dallas, TX, USA. \href{https://arxiv.org/abs/2508.04516}{arXiv:2508.04516}

  \item \textbf{Ishraq Tashdid}, Tasnuva Farheen, Sazadur Rahman,
    \textit{``SAFE-SiP: Secure Authentication Framework for System-in-Package Using
    Multi-party Computation,''} \textbf{GLSVLSI 2025} (ACM Great Lakes Symposium on VLSI),
    New Orleans, LA, USA. \href{https://arxiv.org/abs/2505.09002}{arXiv:2505.09002}

  \item \textbf{Ishraq Tashdid}, Tasnuva Farheen, Sazadur Rahman,
    \textit{``AuthenTree: A Scalable MPC-Based Distributed Trust Architecture for
    Chiplet-based Heterogeneous Systems,''} \textbf{PAINE 2025} (IEEE Physical Assurance
    and Inspection of Electronics), Denver, CO, USA.
    \href{https://arxiv.org/abs/2508.13033}{arXiv:2508.13033}

  \item \textbf{Ishraq Tashdid}, Valentina Terry, Jordan Merkel, Tasnuva Farheen, Sazadur Rahman,
    \textit{``BeyondPPA: Human-Inspired Reinforcement Learning for Post-Route
    Reliability-Aware Macro Placement,''} \textbf{MLCAD 2025} (ACM/IEEE Symposium on
    Machine Learning for CAD), Santa Cruz, CA, USA.
    \uline{\textbf{Best Paper Nominee.}}
    \href{https://doi.org/10.1109/MLCAD65511.2025.11189164}{DOI:10.1109/MLCAD65511.2025.11189164}

  \item \textbf{Book Chapter (invited):} \emph{LLMs for SoC Design and Security} ---
    with collaborators at the University of Florida, \textit{in preparation}, 2025--2026.

\end{itemize}


% ─── INVITED TALKS ───────────────────────────────────────────────────────────
\section*{Invited Talks \& Presentations}

\textbf{Intel AI Forum} --- Intel Corporation, Folsom, CA \hfill \textit{Feb.\ 2026}\\
Invited (non-submitted) talk on \emph{ATLAS}: LLM-driven threat modeling and formal security
verification for SoC security. Presented to Intel's AI research community of engineers and
researchers; selected by Intel organizers from ongoing academic collaborations. Directly led
to selection for Intel's competitive XPU Security Research Internship.

\vspace{3pt}
\textbf{4th Annual Florida Semiconductor Summit (FSI)} --- Rosen Shingle Creek, Orlando, FL
\hfill \textit{Feb.\ 2026}\\
Invited presentation on hardware security and AI-assisted verification for Florida's
semiconductor manufacturing ecosystem. Theme: \emph{Semiconductor Manufacturing in Florida:
Power. Progress. Possibilities.} Organized by Florida Semiconductor Institute.

\vspace{3pt}
\textbf{Open-Source EDA Birds-of-a-Feather} --- 63rd Design Automation Conference (DAC 2026),
Long Beach Convention Center, Long Beach, CA \hfill \textit{Jul.\ 2026}\\
Invited rapid-fire presentation, \emph{``Security Challenges and Opportunities in the
Open-Source EDA Ecosystem,''} at the Open-Source EDA, Data and Collaboration Summit --- a
recurring DAC Birds-of-a-Feather forum (running since 2018) organized by UC San Diego,
Google, and Precision Innovations, bringing together academic researchers, industry
practitioners, and open-source EDA maintainers.

% ─── TEACHING & MENTORING ────────────────────────────────────────────────────
\section*{Teaching \& Mentoring}

\textbf{Graduate Teaching Assistant} --- UCF ECE \hfill \textit{Spring 2024; Summer \& Fall 2025}
\begin{itemize}
  \item Courses: Engineering Analysis and Computation; Linear Systems II;
  Electronics II Lab; Digital Systems Lab. Responsibilities included office hours,
  grading, and lab instruction.
\end{itemize}

\vspace{2pt}
\textbf{Undergraduate Research Mentor} \hfill \textit{Aug.\ 2024--Present}
\begin{itemize}
  \item Mentored two undergraduate researchers in the UCF hardware security lab.
  Both students secured \textbf{summer 2025 research internships at AMD and Northrop Grumman},
  respectively, and co-authored the \textbf{MLCAD 2025 Best Paper Nominee}.
  Northrop Grumman placement reflects direct contribution to the US defense contractor
  engineering pipeline.
\end{itemize}


% ─── PROFESSIONAL SERVICE ────────────────────────────────────────────────────
\section*{Professional Service}

\textbf{Program Committee}, AAAI-26 Fall Symposium Series (FSS-26)\\
\textbf{External Reviewer}, Design Automation Conference (DAC 2026 and DAC 2025)\\
\textbf{External Reviewer}, IEEE International Conference on Computer Design (ICCD 2025)\\
\textbf{External Reviewer}, IEEE International Symposium on Hardware Oriented Security and Trust (HOST 2026)\\
\textbf{External Reviewer}, IEEE Design \& Test (journal)\\
Serving on a program committee and invited by program committees and editorial boards to
review submissions spanning EDA, hardware security, computer design, and artificial
intelligence --- across flagship conferences, an AAAI symposium series, and an archival
IEEE transactions-family journal. These invitations reflect recognition by field leaders
as a qualified independent expert.\\

% ─── PATENTS ─────────────────────────────────────────────────────────────────
\section*{Patents}

\textbf{Provisional Patent Application} \hfill \textit{2026}\\
Filed on the \emph{ATLAS} LLM-driven threat-modeling and automated security-property
generation framework (DAC 2026), with Dr.\ Sazadur Rahman and Kimia Tasnia. \emph{(Provisional/preliminary
stage.)}

% ─── GRANTS & PROPOSALS ──────────────────────────────────────────────────────
\section*{Grants \& Proposals}

\textbf{Research Contributor} \hfill \textit{2026}\\
\textbf{DOE Genesis Mission: Transforming Science and Energy with AI}
(RFA DE-FOA-0003612), PI: Dr.\ Sazadur Rahman, in collaboration with \textbf{Silvaco Inc.}
\emph{(Submitted; pending review.)}\\
Contributed to proposal development and technical research direction under DOE's
\$293M national AI initiative, targeting the semiconductors and microelectronics
challenge area. Technical content builds directly on ATLAS and related security-verification
publications. Reflects early-career engagement at the grant-writing and research
leadership level in a White House-led federal program.

\vspace{4pt}
\textbf{Research Contributor} \hfill \textit{2025}\\
\textbf{ASCEND} --- NSF Secure and Trustworthy Cyberspace (SaTC 2.0) Program,
PI: Dr.\ Sazadur Rahman. \emph{(Submitted; pending review.)}\\
Contributed to the proposal's technical research direction on security and trust for
next-generation heterogeneous chiplet integration and advanced packaging, extending
authentication and trust-architecture methodology from InterPUF, SAFE-SiP, and AuthenTree.

\vspace{4pt}
\textbf{Research Contributor} \hfill \textit{2025}\\
\textbf{SecureOpenROAD} --- NSF Safety, Security, and Privacy in Open Source Ecosystems
(Safe-OSE, NSF 24-608), PI: Dr.\ Sazadur Rahman, in collaboration with the OpenROAD
Initiative (ORI). \emph{(Submitted.)}\\
Contributed to the proposal's technical research direction on strengthening security assurance
within the OpenROAD open-source chip-implementation ecosystem, building on research first
presented at the DAC 2026 Open-Source EDA Birds-of-a-Feather session.

% ─── HONORS & AWARDS ────────────────────────────────────────────────────────
\section*{Honors, Awards \& Fellowships}

\begin{itemize}[leftmargin=*]
  \item \textbf{UCF Trustees Doctoral Fellowship} (2024--2028) --- 4-year fellowship (\$25{,}000/year)
  awarded to newly enrolled doctoral students with outstanding records. Among the most
  selective doctoral awards at UCF; recipients are identified for exceptional academic
  and research potential at the time of admission.

  \item \textbf{UCF ECE Outstanding Student Researcher Award} (May 2026) --- Departmental
  recognition for outstanding contributions to research during FY 2025--26, awarded by the
  UCF Department of Electrical and Computer Engineering.

  \item \textbf{DAC Young Fellow} --- 63rd Design Automation Conference (DAC 2026), Long Beach,
  CA. Competitive recognition selecting outstanding early-career researchers at the premier
  international conference on electronic design automation.

  \item \textbf{Future Faculty Laureate (FFL)} --- UCF ECE Department (Fall 2024--Present)\\
  Competitive program identifying Ph.D.\ students with exceptional potential for academic
  faculty careers. Participants receive mentorship, teaching development, and professional
  preparation for faculty positions at research universities.

  \item \textbf{Best Paper Award Nominee} --- ACM/IEEE MLCAD 2025\\
  Nominated by program committee among top papers at the Symposium on Machine Learning for CAD.

  \item \textbf{NSF Travel Grant} --- IEEE International Symposium on Hardware Oriented
  Security and Trust (HOST 2025). Federal recognition of research merit; awarded competitively
  to support participation in the premier hardware security symposium.

\end{itemize}


% ─── SKILLS ──────────────────────────────────────────────────────────────────
\section*{Technical Skills}
\begin{tabular}{@{}p{4.4cm}p{12.0cm}@{}}
\textbf{Security \& Analysis:}   & Threat modeling (CWE/CVE/CAPEC), LLM-assisted vulnerability
                                    reasoning, assertion-based verification, formal verification
                                    (JasperGold), static analysis, hardware IP protection \\
\textbf{Programming:}            & Python, SystemVerilog, Verilog, C, MATLAB, Assembly (RISC-V) \\
\textbf{Verification:}           & UVM, cocotb, pytest, ABV; RISC-V ISA \& system-level
                                    compliance; AXI/APB bus protocols \\
\textbf{EDA \& Simulation:}      & Cadence Innovus / Genus / JasperGold, ModelSim,
                                    Quartus Prime, Yosys, PyVerilog \\
\textbf{ML / AI:}                & PyTorch, reinforcement learning (PPO/DDPG),
                                    LLM prompting \& fine-tuning, RAG pipelines \\
\end{tabular}

\vspace{10pt}
\begin{center}
{\footnotesize \textcolor{gray}{Last updated: September 2026}}
\end{center}

\end{document}
